\begin{table}[!ht]
\centering
\scriptsize
\setlength{\tabcolsep}{3pt}
\caption{Propensity-model sensitivity on the SYM-NL and SYM-MU designs at $n=500$ and $n=1000$ over fifty paired replications. Entries are end-to-end reference-TCATE Monte Carlo RMSEs with Monte Carlo standard errors in parentheses; lower is better. The columns replace the default logistic propensity with a random-forest factory, a gradient-boosting factory, and the oracle propensity, which is a diagnostic and not a feasible competitor. The outcome learner, grid, and particle budget are held fixed.}
\label{tab:sensitivity-propensity}
\begin{tabular}{lrrrr}
\toprule
DGP & Logistic & Random forest & Gradient boosting & Oracle \\
\midrule
\multicolumn{5}{l}{\textit{Panel A: $n=500$}} \\
\midrule
SYM-NL & 0.0562(0.0059) & 0.0575(0.0060) & 0.0575(0.0060) & 0.0575(0.0060) \\
SYM-MU & 0.0487(0.0054) & 0.0487(0.0045) & 0.0487(0.0045) & 0.0487(0.0045) \\
\multicolumn{5}{l}{\textit{Panel B: $n=1000$}} \\
\midrule
SYM-NL & 0.0420(0.0057) & 0.0427(0.0055) & 0.0427(0.0055) & 0.0427(0.0055) \\
SYM-MU & 0.0361(0.0045) & 0.0348(0.0042) & 0.0348(0.0042) & 0.0348(0.0042) \\
\bottomrule
\end{tabular}
\vspace{2pt}
{\footnotesize \textit{Note:} the confirmatory protocol reports the common cross-fitted AIPW view, whose propensity layer is fixed across columns, so the factory choice has little effect on the tabulated errors. The estimator's internal reference-estimate influence-function SE, a calibration diagnostic, is logistic 0.021, random forest 0.022, gradient boosting 0.079, and oracle 0.022; the gradient-boosting factory is roughly three to four times less stable than the alternatives, consistent with the archive's miscalibration diagnosis.}
\end{table}
